% !TeX program = xelatex
\documentclass[11pt,a4paper]{article}
\usepackage{fontspec}
\usepackage[spanish, es-noshorthands]{babel}
\usepackage{amsmath}
\usepackage{graphicx}
\usepackage{booktabs}
\usepackage{geometry}
\geometry{margin=2.5cm}
\usepackage{tikz}
\usetikzlibrary{arrows.meta, positioning}
\usepackage{caption}
\usepackage{hyperref}

\title{Diseño experimental del batch instrumentado para la\\
bioconversión de residuos con \emph{Hermetia illucens}:\\
reactor, instrumentación, control y protocolos}
\author{J.~J.~Leal \and Grupo de investigación}
\date{Versión 1.0 --- \today}

\begin{document}
\maketitle
\tableofcontents
\newpage

\section{Objetivo y alcance}

Construir y operar un batch de bioconversión con \emph{Hermetia illucens}
totalmente instrumentado, capaz de: (i) medir emisiones de GEI
(CO$_2$, CH$_4$) con separación de fuentes (larvaria vs.\ microbiana),
(ii) controlar temperatura, humedad y ventilación de la cama,
(iii) alimentar y validar el gemelo digital del proceso, y
(iv) generar los datos del artículo Q1.
El diseño parte de los resultados ya obtenidos con la cámara
tipo panera de 11{,}9~L y los modelos 0D/2D del gemelo.

\section{Reactor: panera modificada}

Recipiente plástico tronco-piramidal con tapa (medido):
base $31\times20$~cm, boca $35\times25$~cm, altura 16~cm;
$V_\text{total}=11{,}9$~L, $V_\text{aire}\approx11{,}4$~L con cama.

Modificaciones obligatorias:
\begin{enumerate}
  \item \textbf{Placa difusora de aireación} bajo la cama
        (acrílico perforado 3~mm, orificios cada 2~cm, plenum de 1~cm).
  \item \textbf{Puerto de muestreo con septo} en la tapa
        (septo de silicona/PTFE, para jeringa hermética).
  \item \textbf{Válvula de modo dual} en la salida:
        cerrada (medición de flujo tipo cámara estática) /
        abierta (operación continua ventilada).
  \item Pasacables herméticos (prensas PG7 + silicona) para sensorería.
\end{enumerate}

\textbf{Prueba de fugas (P-01):} inyectar CO$_2$ hasta 3\,000~ppm,
cerrar, registrar 1~h. Criterio de aceptación: decaimiento $<5$\,\%/h.

\section{Sistema de aireación forzada de fondo}

\textbf{Justificación (resultados propios):} los modelos 2D y 0D
mostraron que ventilar la cabecera no controla ni el CO$_2$ ni la
temperatura \emph{dentro} de la cama: la difusión es el cuello de
botella. La aireación debe entrar \emph{por debajo}.

\subsection{Caudales de diseño (calculados de datos reales)}
Flujo máximo medido: 294~mL\,CO$_2$/h con 750 larvas (D1, agitadas).
Escalando al batch nominal (2\,000 larvas): $q_\text{max}\approx0{,}8$~L/h.
Para mantener CO$_2\leq0{,}45$\,\% bajo el tope del sensor:
\[
  \text{ACH} \geq \frac{100\,q_\text{max}}{V_\text{aire}\cdot 0{,}45}
  \approx 15~\text{recambios/h}.
\]
Rango de diseño: \textbf{0{,}5--20 ACH}, es decir
$Q = 0{,}1$--$3{,}8$~L/min. Componentes:
\begin{itemize}
  \item Bomba de diafragma 5~L/min (acuario, 2 salidas).
  \item Control de caudal: rotámetro 0{,}1--5~L/min con válvula
        de aguja (básico) o controlador de flujo másico (ideal).
  \item Humidificador de burbujeo en la entrada (frasco lavador)
        para no secar la cama.
  \item Filtro de entrada (HEPA de acuario) contra contaminación biológica.
\end{itemize}
\textbf{Ubicación inlet:} plenum bajo la cama (distribución uniforme).
\textbf{Ubicación outlet:} esquina superior de la tapa, diagonalmente
opuesta al punto de medición, con válvula solenoide de modo dual.

\section{Control térmico}

Baño termostático: la panera se instala dentro de un contenedor mayor
con agua a $30{,}0\pm0{,}5$~°C (calentador de acuario de 100~W con
termostato + minibomba de circulación). La panera se apoya sobre
pilares sellados que atraviesan el fondo del baño y transmiten su peso
a las celdas de carga (sección~\ref{sec:instrumentacion}).
Justificación: óptimo larvario 25--30~°C; el baño fija la frontera
térmica y la aireación hace el control fino de la cama.

\section{Control de humedad}

\begin{itemize}
  \item \textbf{Cama:} 60--70\,\% (óptimo literatura). Medición:
        sensor capacitivo enterrado + celda de carga (balance de agua).
        Corrección: aspersión medida por jeringa sobre la superficie,
        nunca por encima de 70\,\% (riesgo anaerobio).
  \item \textbf{Cabecera:} HR 60--70\,\%, controlada indirectamente
        por el humidificador de entrada y el caudal.
\end{itemize}

\section{Instrumentación}
\label{sec:instrumentacion}

\begin{table}[h]
\centering\small
\caption{Instrumentación del batch. Ubicaciones en
figura~\ref{fig:reactor}.}
\begin{tabular}{@{}llllll@{}}
\toprule
Variable & Sensor & Rango & Exactitud & Ubicación & Frec. \\
\midrule
CO$_2$ & Winsen MH-410D (v.~5\,\%vol) & 0--50\,000~ppm & $\pm20$~ppm (validado GC) & cabecera centro & 30~s \\
CH$_4$ & Winsen MH-441D & 0--5\,\%vol & $\pm3$\,\% FS; validar COVs vs GC & cabecera centro & 30~s \\
O$_2$  & SST LuminOx LOX-02 (óptico) & 0--25\,\% & $\pm2$\,\% FS; T90$<$30~s & cabecera centro & 30~s \\
T/HR aire & Sensirion SHT45 & $-40$--125~°C, 0--100\,\% & $\pm0{,}1$~°C, $\pm1{,}5$\,\%HR & cabecera centro & 30~s \\
T cama & DS18B20 estanco ($\times3$) & $-55$--125~°C & $\pm0{,}5$~°C & prof.\ 1, 2{,}5 y 4~cm & 30~s \\
Humedad cama & capacitivo v2.0 & 0--100\,\% & calibrar gravimétrico & centro de cama & 30~s \\
Masa total & 4$\times$ celda 50~kg + HX711 & 0--200~kg & $\pm10$~g & bajo la panera & 30~s \\
Caudal & rotámetro 0{,}1--5~L/min & --- & $\pm5$\,\% & línea de entrada & manual \\
\bottomrule
\end{tabular}
\end{table}

\textbf{Nota crítica CO$_2$:} la versión de 5\,000~ppm queda
descartada para el batch (saturó en campaña anterior); se requiere
la versión 5\,\%vol. El O$_2$ permite calcular el cociente
respiratorio RQ = CO$_2$/O$_2$, variable metabólica clave.

\section{Diseño de la cama y especificaciones del batch}
\label{sec:batch}

Basado en la literatura de cría (óptimos: T 25--30~°C, humedad de
sustrato 60--70\,\%, HR 60--70\,\%, densidad 2{,}5--7
larvas/cm$^2$, profundidad 1--5~cm, oscuridad):

\begin{table}[h]
\centering\small
\caption{Especificaciones del batch nominal.}
\begin{tabular}{@{}lll@{}}
\toprule
Parámetro & Valor nominal & Base \\
\midrule
Larvas & 2\,000 (pesadas, $\pm$10) & 2{,}7/cm$^2$ sobre 742~cm$^2$ \\
Edad siembra & 4--6 días (neonatas tardías) & manejo y viabilidad \\
Alimento & 100~mg/larva/día (base húmeda) & rango 90--175~mg/larva/día \\
Cargas & días 0, 3 y 6 (tercios) & evita compactación anaerobia \\
Profundidad cama & 3--4{,}5~cm & máximo 5~cm \\
Humedad sustrato & 65--70\,\% & óptimo digestión \\
Temperatura & $27\pm1$~°C baño; cama $\leq33$~°C & bienestar larvario \\
HR cabecera & 60--70\,\% & --- \\
Iluminación & oscuridad & instinto de enterramiento \\
Duración & 12--14 días o señal de cosecha & --- \\
\bottomrule
\end{tabular}
\end{table}

\section{Diseño factorial y controles}

Factores: tipo de alimento (2) $\times$ aireación (sin / cabecera /
fondo) $\times$ día larvario (D1--D14, medición diaria).
Controles por batch: (a) alimento solo por tipo (microbiano),
(b) cámara vacía (fondo de sensores), (c) cámara cerrada con larvas
sin alimento (ventana corta, separación larvaria pura).
Réplicas: $n\geq3$ por condición. \textbf{Pesar larvas en siembra y
cosecha, y cada 2 días en muestreo destructivo paralelo si es viable.}

\section{Protocolos operativos}

\begin{description}
  \item[P-01 Fugas:] cámara vacía, CO$_2$ a 3\,000~ppm, 1~h;
        aceptación $<5$\,\%/h.
  \item[P-02 Validación de sensores:] (i) test de etanol para
        MH-441D (paño con alcohol junto a entrada, cámara vacía);
        (ii) cero y span con gas patrón; (iii) registro de deriva.
  \item[P-03 Arranque:] cama preparada 24~h antes (fermentación
        controlada), siembra, peso inicial, inicio de registro.
  \item[P-04 Operación diaria:] ventana de cámara cerrada de
        15~min (cerrar válvula, aireación en cero) con criterio de
        parada CO$_2<4\,000$~ppm; en el minuto 12, muestreo por
        septo para GC (2$\times$20~mL, CH$_4$, CO$_2$, etanol/COVs).
  \item[P-05 Cosecha:] por señal (sección~\ref{sec:cosecha}) o día
        14; peso de larvas, frass, supervivencia.
  \item[P-06 Limpieza:] lavado, calibración de sensores, respaldo
        de datos al repositorio.
\end{description}

\section{Señales de cosecha (hipótesis a probar)}
\label{sec:cosecha}

Tres candidatas, a validar contra biomasa real cosechada:
(i) pico de CH$_4$/COVs (transición anaerobia del alimento),
(ii) caída del CO$_2$ tras el pico (metabolismo cayendo por
agotamiento de alimento), (iii) aplanamiento de la pendiente de masa
(celda de carga). El diseño mide las tres de forma continua.

\section{Control y gemelo digital}

Variables manipuladas: caudal de aireación, válvula de modo,
calefactor del baño, eventos de rehidratación.
Leyes iniciales (ya simuladas en el gemelo): ACH adaptativo para
CO$_2<0{,}45$\,\%, ACH máximo si $T_\text{cama}>33$~°C.
Camino siguiente: calibración inversa con PINN y control predictivo
(MPC) informado por el gemelo.

\section{Adquisición de datos}

ESP32/Raspberry Pi; lectura cada 30~s (mismo formato TSV actual:
\texttt{hh:mm:ss}, valor); CSV diario con marca de eventos
(ventana cerrada, muestreo GC, recarga de alimento); respaldo
automático al repositorio \texttt{tesis-BSF/datos/}.

\section{Diagrama del reactor}

\begin{figure}[h]
\centering
\begin{tikzpicture}[font=\footnotesize, >=Stealth]

% --- Bano termico externo ---
\draw[fill=cyan!8] (0,0) rectangle (12.4,5.6);
\node[anchor=west, text=teal] at (0.25,5.25) {baño térmico 30 $\pm$ 0.5 °C};

% --- Panera ---
\draw[thick, fill=white] (1.6,1.1) -- (1.9,4.3) -- (10.7,4.3) -- (11.0,1.1) -- cycle;
\draw[thick, fill=red!15] (1.75,4.3) rectangle (10.85,4.7); % tapa

% --- Cama de sustrato ---
\draw[fill=brown!35] (1.73,1.45) rectangle (10.9,2.25);
\node at (6.3,1.85) {cama: sustrato + larvas (3--4.5 cm)};

% --- Placa difusora / plenum ---
\draw[fill=gray!45] (1.73,1.1) rectangle (10.9,1.35);
\node[anchor=west] at (11.15,1.22) {\parbox{3.1cm}{placa difusora\\+ plenum (aireación)}};

% --- Pilares de peso ---
\draw[fill=black!60] (3.0,0.55) rectangle (3.3,1.1);
\draw[fill=black!60] (9.3,0.55) rectangle (9.6,1.1);
\draw[fill=orange!30] (2.7,0.15) rectangle (3.6,0.55);
\draw[fill=orange!30] (9.0,0.15) rectangle (9.9,0.55);
\node[anchor=north, align=center] at (3.15,0.12) {celda de carga};
\node[anchor=north, align=center] at (9.45,0.12) {celda de carga};

% --- Linea de entrada de aire ---
\node[draw, fill=green!15, rounded corners, align=center] (bomba) at (-2.2,1.1) {bomba\\5 L/min};
\node[draw, fill=green!15, rounded corners, align=center, right=0.45cm of bomba] (filtro) {filtro};
\node[draw, fill=green!15, rounded corners, align=center, right=0.45cm of filtro] (hum) {humidif.};
\node[draw, fill=green!15, rounded corners, align=center, right=0.45cm of hum] (flujo) {caudal\\0.1--5 L/min};
\draw[->, thick, green!50!black] (bomba) -- (filtro);
\draw[->, thick, green!50!black] (filtro) -- (hum);
\draw[->, thick, green!50!black] (flujo.east) -- ++(0.35,0) |- (6.3,1.1);
\draw[->, thick, green!50!black] (hum) -- (flujo);

% --- Salida con valvula de modo dual ---
\node[draw, fill=violet!15, rounded corners, align=center] (valv) at (12.6,4.5) {válvula\\modo dual};
\draw[->, thick, violet] (10.85,4.5) -- (valv.west);
\draw[->, thick, violet] (valv.east) -- ++(0.9,0) node[midway, above]{salida};

% --- Sensores de cabecera (cuelgan de la tapa) ---
\foreach \x/\lab in {3.4/{CO$_2$}, 5.0/{CH$_4$}, 6.6/{O$_2$}, 8.2/{T/HR}}{
  \draw[fill=yellow!25] (\x-0.35,3.55) rectangle (\x+0.35,4.3);
  \node at (\x,3.92) {\lab};
}
\node[anchor=west, align=left] at (11.15,3.1) {sensores cabecera\\(centro, lejos de salida)};

% --- Septo para GC ---
\draw[fill=magenta!30] (9.4,4.3) rectangle (10.0,4.55);
\node[anchor=south west, align=left] at (9.2,4.8) {septo (muestreo GC)};

% --- Sondas T en la cama ---
\foreach \y in {1.6,1.85,2.1}{
  \draw[fill=red!60] (2.3,\y) circle (0.06);
}
\node[anchor=east, align=right] at (2.2,1.85) {3$\times$ T cama\\(1, 2.5, 4 cm)};

% --- Unidad de control ---
\node[draw, fill=blue!10, rounded corners, align=center] (ctrl) at (-2.2,4.3) {ESP32 / RPi\\control + datos};
\draw[dashed, blue!60] (ctrl) -- (valv.north west);
\draw[dashed, blue!60] (ctrl) |- (bomba.north);
\draw[dashed, blue!60] (ctrl.east) -- ++(0.3,0) |- (5.0,4.75);

\end{tikzpicture}

\caption{Sección transversal esquemática del reactor instrumentado.}
\label{fig:reactor}
\end{figure}

\begin{thebibliography}{9}\small
\bibitem{mertenat2019} Mertenat, A., Diener, S., Zurbrügg, C. (2019).
Black Soldier Fly biowaste treatment -- assessment of global warming
potential. \emph{Waste Management} 84, 173--181.
\bibitem{parodi2020} Parodi, A.\ et al. (2020). Bioconversión de
residuos orgánicos con H.~illucens y balance de GEI.
\emph{Journal of Cleaner Production}. [TODO: completar desde Zotero]
\bibitem{eriksson2022} Eriksson et al. (2022). Respiración y CO$_2$
en larvas de H.~illucens. [TODO: completar desde Zotero]
\bibitem{welfare2024} Anon. (2024). Black Soldier Fly Larvae's
Optimal Feed Intake and Rearing Density: A Welfare Perspective
(Part~II). \emph{Animals} (MDPI). [TODO: completar desde Zotero]
\bibitem{feeding2026} Anon. (2026). Feeding Strategies for Optimizing
BSF Larval Production. \emph{Sustainability} (MDPI). [TODO: completar]
\bibitem{mh410d} Winsen. Datasheet MH-410D NDIR CO$_2$ sensor.
\bibitem{mh441d} Winsen. Datasheet MH-441D NDIR CH$_4$ sensor.
\bibitem{lox02} SST Sensing. LuminOx LOX-02 optical O$_2$ sensor,
DS-0030.
\bibitem{sht45} Sensirion. Datasheet SHT4x (SHT45).
\end{thebibliography}

\end{document}
